\documentclass[11pt]{article}
\usepackage[T1]{fontenc}
\usepackage[utf8]{inputenc}
\usepackage[margin=0.92in]{geometry}
\usepackage{amsmath,amssymb,amsthm}
\usepackage{graphicx}
\usepackage{microtype}
\usepackage[hidelinks]{hyperref}
\usepackage{enumitem}

\newtheorem{theorem}{Theorem}
\newtheorem{lemma}{Lemma}
\newcommand{\K}{\mathbb K}
\newcommand{\JNP}{\mathrm{JNP}}
\newcommand{\Cp}{C_p}
\newcommand{\bbeta}{\beta^{\!*}}

\title{A finite-slice reduction for Josefson--Nissenzweig witnesses in $C_p(X,E)$}
\author{Candidate partial result for arXiv:2107.03211, Problem 4.8}
\date{19 July 2026}

\begin{document}
\maketitle

\noindent\textbf{Status.}
Unconditional partial result, likely valid, pending human review.  The
unrestricted problem remains open after the bounded literature check recorded
below.

\section{Source problem}

Bargetz, K\k{a}kol, and Sobota ask the following in Problem 4.8 on page 7 of
arXiv:2107.03211v3 \cite{BKS}.

\begin{quote}
Let $X$ be a Tychonoff space and $E$ a locally convex space.  Assume that
$C_p(X,E)$ has the Josefson--Nissenzweig property.  Does it follow that
$C_p(X)$ has the Josefson--Nissenzweig property or $E$ has the
Josefson--Nissenzweig property?
\end{quote}

\begin{figure}[h]
  \centering
  \includegraphics[width=\linewidth]{figures/open_problem_crop.png}
  \caption{Problem 4.8 and its immediate context in the source paper.}
\end{figure}

\section{Definitions and canonical form}

For a locally convex space $G$, a set $B\subseteq G$ is
\emph{barrel-bounded} if every barrel in $G$ absorbs $B$.  The topology
$\beta^*(G',G)$ is uniform convergence on barrel-bounded subsets of $G$.
The space $G$ has the JNP when there is a sequence in $G'$ which is null for
$\sigma(G',G)$ but not null for $\beta^*(G',G)$; see
\cite{BG,BKS}.

Set $H=C_p(X,E)$.  The source paper recalls the algebraic identification
\[
  H'=L_p(X)\otimes E'.
\]
After grouping equal evaluation points, every $\Phi\in H'$ has the form
\begin{equation}\label{eq:grouped}
  \Phi(g)=\sum_{x\in A}\eta_x(g(x)),
  \qquad A\subseteq X\text{ finite},\quad \eta_x\in E'.
\end{equation}
The nonzero coefficients in this form are unique.  Indeed, a Tychonoff space
admits continuous interpolation on each finite subset: for distinct
$x_1,\dots,x_m$ there are $h_j\in C(X)$ with
$h_j(x_k)=\delta_{jk}$.

\section{Proof intuition}

A finite slice lets us separate the tensor variables.  If only finitely many
points of $X$ occur, finite interpolation extracts a weak-star null sequence
in $E'$ at each point.  If only finitely many coefficient directions in $E'$
occur, a biorthogonal family in $E$ extracts finitely many weak-star null
sequences in $L_p(X)$.  Evaluation and scalarization are continuous linear
maps, so they carry a barrel-bounded witnessing set to barrel-bounded sets.
If every extracted component converged uniformly there, their finite sum
would also converge uniformly, contradicting the original JNP witness.

Thus neither kind of one-sided finite confinement can produce a genuinely
mixed counterexample to Problem 4.8.  The remaining case is precisely the
two-sided escape regime: infinitely many evaluation points and an
infinite-dimensional coefficient span.

\section{Finite-slice theorem}

\begin{lemma}[barrel-bounded images]\label{lem:image}
If $T:G\to F$ is continuous linear and $B\subseteq G$ is barrel-bounded,
then $T(B)$ is barrel-bounded in $F$.
\end{lemma}

\begin{proof}
For every barrel $V\subseteq F$, the inverse image $T^{-1}(V)$ is closed,
absolutely convex, and absorbing in $G$, hence is a barrel.  It absorbs $B$,
so $V$ absorbs $T(B)$.
\end{proof}

\begin{theorem}[finite-slice reduction]\label{thm:main}
Let $X$ be Tychonoff, let $E$ be locally convex, and put $H=C_p(X,E)$.
Suppose $(\Phi_n)\subseteq H'$ is weak-star null but not
$\beta^*(H',H)$-null.  Write
\[
  \Phi_n(g)=\sum_{x\in A_n}\eta_{n,x}(g(x))
\]
in grouped form.  Then:
\begin{enumerate}[label=\textup{(\roman*)}]
  \item if $\bigcup_n A_n$ is finite, then $E$ has the JNP;
  \item if
  $\operatorname{span}\{\eta_{n,x}:n\in\mathbb N,\ x\in A_n\}$ is
  finite-dimensional, then $C_p(X)$ has the JNP.
\end{enumerate}
Consequently, if neither $C_p(X)$ nor $E$ has the JNP, every JNP witness in
$C_p(X,E)'$ must escape every finite evaluation slice and every
finite-dimensional coefficient slice.
\end{theorem}

\begin{proof}
Because $(\Phi_n)$ is not $\beta^*$-null, there is a barrel-bounded set
$B\subseteq H$ on which uniform convergence fails.

For (i), let $A=\bigcup_n A_n$ and put $\eta_{n,x}=0$ when
$x\notin A_n$.  Fix $x\in A$.  Choose $h_x\in C(X)$ which is $1$ at $x$
and $0$ on $A\setminus\{x\}$.  For $e\in E$, the function
$g_{x,e}(t)=h_x(t)e$ belongs to $H$, and
\[
  \eta_{n,x}(e)=\Phi_n(g_{x,e})\longrightarrow 0.
\]
Thus $(\eta_{n,x})_n$ is weak-star null in $E'$ for every $x\in A$.
By Lemma \ref{lem:image}, the set
$B_x=\{g(x):g\in B\}$ is barrel-bounded in $E$.  If every
$(\eta_{n,x})_n$ were $\beta^*(E',E)$-null, then
\[
  \sup_{g\in B}|\Phi_n(g)|
  \leq \sum_{x\in A}\sup_{e\in B_x}|\eta_{n,x}(e)|
  \longrightarrow 0,
\]
contrary to the choice of $B$.  Hence one component witnesses the JNP of
$E$.

For (ii), let $F\subseteq E'$ be the finite-dimensional span of all
coefficients and choose a basis $e'_1,\dots,e'_m$ of $F$.  There are
$e_1,\dots,e_m\in E$ with $e'_j(e_k)=\delta_{jk}$.  To see this, the map
$E\to\K^m$, $e\mapsto(e'_1(e),\dots,e'_m(e))$, is onto; otherwise a
nonzero functional on $\K^m$ annihilating its range would give a linear
dependence among the $e'_j$.

Write
\[
  \eta_{n,x}=\sum_{j=1}^m a_{n,x,j}e'_j,
  \qquad
  \mu_{n,j}=\sum_{x\in A_n}a_{n,x,j}\delta_x\in L_p(X).
\]
For $f\in C(X)$, apply $\Phi_n$ to $t\mapsto f(t)e_j$ to obtain
$\mu_{n,j}(f)\to0$.  Hence $(\mu_{n,j})_n$ is weak-star null in
$C_p(X)'$.

The scalarization map $S_j:H\to C_p(X)$,
$S_j(g)=e'_j\circ g$, is continuous linear.  Thus
$D_j=S_j(B)$ is barrel-bounded by Lemma \ref{lem:image}.  If every
$(\mu_{n,j})_n$ were $\beta^*(C_p(X)',C_p(X))$-null, then
\[
  \sup_{g\in B}|\Phi_n(g)|
  \leq \sum_{j=1}^m\sup_{f\in D_j}|\mu_{n,j}(f)|
  \longrightarrow0,
\]
again a contradiction.  One component therefore witnesses the JNP of
$C_p(X)$.
\end{proof}

\section{Scope, novelty, and verification}

Theorem \ref{thm:main} does not settle Problem 4.8: a weak-star null tensor
sequence can have both unbounded evaluation support and
infinite-dimensional coefficient span.  It shows that any counterexample must
use both phenomena essentially, rather than merely embedding a witness from
one factor in a disguised finite slice.

The proof is algebraic-topological and uses no computation.  The main audit
points are Lemma \ref{lem:image}, finite interpolation on $X$, and the finite
sum estimate.  All are explicit above.

A bounded novelty search covered the run's lightweight indexes and exact
arXiv queries using the source title, Problem 4.8, $C_p(X,E)$, JNP,
JN-sequences, $L_p(X)\otimes E'$, finite support, and finite-dimensional
coefficient span.  The background paper arXiv:2003.06764 \cite{BG} and the
latest source version were checked.  No exact finite-slice result or later
solution of Problem 4.8 was found.  Novelty is plausible but not certified.

\paragraph{Human-review recommendation.}
Review as a likely valid partial reduction.  The priority is to verify the
barrel-bounded image step and then decide whether known tensor-sequence
compactness principles can force one of the two finite-slice alternatives.

\begin{thebibliography}{9}
\bibitem{BKS}
C. Bargetz, J. K\k{a}kol, and D. Sobota,
\emph{On complemented copies of the space $c_0$ in spaces $C_p(X,E)$},
arXiv:2107.03211v3, 2023.

\bibitem{BG}
T. Banakh and S. Gabriyelyan,
\emph{The Josefson--Nissenzweig property for locally convex spaces},
arXiv:2003.06764, 2020.
\end{thebibliography}

\end{document}
